\section{OOL: Onset--Offset Loss}

Ground-truth onset and offset labels are derived from annotations:
$y^{\on}_{c,t}=1$ on the first frame of each event, $y^{\off}_{c,t}=1$ on the
first frame after each event, and zero elsewhere. These labels are extremely
sparse: a 10-second clip has only a handful of boundary frames out of 500.
Plain BCE would be dominated by negatives, so the paper uses a focal loss with
focusing exponent $2$ (the paper writes this as $\alpha=2$; in the focal-loss
literature the exponent is usually $\gamma$):
\begin{equation}
\mathcal{L}_{\on/\off}
=
-\frac{1}{CT}
\sum_{c=1}^{C}\sum_{t=1}^{T}
\begin{cases}
(1-\hat p_{c,t})^{\alpha}\log \hat p_{c,t}, & y_{c,t}=1,\\[3pt]
\hat p_{c,t}^{\alpha}\log(1-\hat p_{c,t}), & y_{c,t}=0.
\end{cases}
\end{equation}
The implementation computes both terms separately:
\begin{equation}
\mathcal{L}_{\mathrm{ool}}
=
\mathcal{L}_{\on}+\mathcal{L}_{\off}.
\end{equation}
For a positive boundary frame with low predicted probability, the factor
$(1-\hat p)^2$ is near one and the loss is large. For an easy negative with low
probability, $\hat p^2$ is tiny and the loss is suppressed. Because RED couples
presence, onset and offset, these losses also shape the presence curve.

